% **** Dokumentacja projektu -- Programowanie obiektowe ****

\documentclass[polish,12pt,twoside,a4paper]{report}
\renewcommand{\baselinestretch}{1.15}

% *************** Definicje stylu dokumentu ***************

% *********************************************************************************
% W pliku tym zdefiniowany jest wygl¹d dokumentu.
% Zmiany tutaj nie s¹ konieczne o ile nie zamierzasz zmieniaæ wygl¹du dokumentu.
% *********************************************************************************

% *************** Za³adowanie pakietów ***************
\usepackage[a4paper,twoside,left=2.0cm,right=1.5cm,top=1.5cm,bottom=1.5cm]{geometry}
\usepackage[T1]{fontenc}
%\usepackage[cp1250]{inputenc}
\usepackage[utf8]{inputenc}
\usepackage[polish]{babel}
\usepackage{amsmath}
\usepackage{amsfonts}
\usepackage{graphicx}
\usepackage{graphics}
\usepackage{times}
\usepackage{indentfirst}%wciecia a nowych akapitach

\selectlanguage{polish}

%szerokoœœ wciêæ
\setlength{\parindent}{1.25cm}

%numeracja stron
\usepackage{fancyhdr}
\pagestyle{fancy}
\fancyhf{} % usun biezace ustawienia pagin
\fancyhead[LE,RO]{ }
\fancyhead[LO]{ }
\fancyhead[RE]{ }
\fancyfoot[CE,CO]{\small\thepage}
\fancyfoot[LO]{ }
\fancyfoot[RE]{ }
\renewcommand{\headrulewidth}{0.0pt}
\renewcommand{\footrulewidth}{0.0pt}
\addtolength{\headheight}{0.0pt} % pionowy odstep na kreske
\fancypagestyle{plain}{%
\fancyhead{} % usun p. górne na stronach pozbawionych
% numeracji (plain)
\renewcommand{\headrulewidth}{0.0pt} % pozioma kreska
\setlength{\headheight}{15pt}
}

% *************** Definicje niektórych kolorów ***************
\usepackage{color}

\definecolor{greenyellow}   {cmyk}{0.15, 0   , 0.69, 0   }
\definecolor{yellow}        {cmyk}{0   , 0   , 1   , 0   }
\definecolor{goldenrod}     {cmyk}{0   , 0.10, 0.84, 0   }
\definecolor{dandelion}     {cmyk}{0   , 0.29, 0.84, 0   }
\definecolor{apricot}       {cmyk}{0   , 0.32, 0.52, 0   }
\definecolor{peach}         {cmyk}{0   , 0.50, 0.70, 0   }
\definecolor{melon}         {cmyk}{0   , 0.46, 0.50, 0   }
\definecolor{yelloworange}  {cmyk}{0   , 0.42, 1   , 0   }
\definecolor{orange}        {cmyk}{0   , 0.61, 0.87, 0   }
\definecolor{burntorange}   {cmyk}{0   , 0.51, 1   , 0   }
\definecolor{bittersweet}   {cmyk}{0   , 0.75, 1   , 0.24}
\definecolor{redorange}     {cmyk}{0   , 0.77, 0.87, 0   }
\definecolor{mahogany}      {cmyk}{0   , 0.85, 0.87, 0.35}
\definecolor{maroon}        {cmyk}{0   , 0.87, 0.68, 0.32}
\definecolor{brickred}      {cmyk}{0   , 0.89, 0.94, 0.28}
\definecolor{red}           {cmyk}{0   , 1   , 1   , 0   }
\definecolor{orangered}     {cmyk}{0   , 1   , 0.50, 0   }
\definecolor{rubinered}     {cmyk}{0   , 1   , 0.13, 0   }
\definecolor{wildstrawberry}{cmyk}{0   , 0.96, 0.39, 0   }
\definecolor{salmon}        {cmyk}{0   , 0.53, 0.38, 0   }
\definecolor{carnationpink} {cmyk}{0   , 0.63, 0   , 0   }
\definecolor{magenta}       {cmyk}{0   , 1   , 0   , 0   }
\definecolor{violetred}     {cmyk}{0   , 0.81, 0   , 0   }
\definecolor{rhodamine}     {cmyk}{0   , 0.82, 0   , 0   }
\definecolor{mulberry}      {cmyk}{0.34, 0.90, 0   , 0.02}
\definecolor{redviolet}     {cmyk}{0.07, 0.90, 0   , 0.34}
\definecolor{fuchsia}       {cmyk}{0.47, 0.91, 0   , 0.08}
\definecolor{lavender}      {cmyk}{0   , 0.48, 0   , 0   }
\definecolor{thistle}       {cmyk}{0.12, 0.59, 0   , 0   }
\definecolor{orchid}        {cmyk}{0.32, 0.64, 0   , 0   }
\definecolor{darkorchid}    {cmyk}{0.40, 0.80, 0.20, 0   }
\definecolor{purple}        {cmyk}{0.45, 0.86, 0   , 0   }
\definecolor{plum}          {cmyk}{0.50, 1   , 0   , 0   }
\definecolor{violet}        {cmyk}{0.79, 0.88, 0   , 0   }
\definecolor{royalpurple}   {cmyk}{0.75, 0.90, 0   , 0   }
\definecolor{blueviolet}    {cmyk}{0.86, 0.91, 0   , 0.04}
\definecolor{periwinkle}    {cmyk}{0.57, 0.55, 0   , 0   }
\definecolor{cadetblue}     {cmyk}{0.62, 0.57, 0.23, 0   }
\definecolor{cornflowerblue}{cmyk}{0.65, 0.13, 0   , 0   }
\definecolor{midnightblue}  {cmyk}{0.98, 0.13, 0   , 0.43}
\definecolor{navyblue}      {cmyk}{0.94, 0.54, 0   , 0   }
\definecolor{royalblue}     {cmyk}{1   , 0.50, 0   , 0   }
\definecolor{blue}          {cmyk}{1   , 1   , 0   , 0   }
\definecolor{cerulean}      {cmyk}{0.94, 0.11, 0   , 0   }
\definecolor{cyan}          {cmyk}{1   , 0   , 0   , 0   }
\definecolor{processblue}   {cmyk}{0.96, 0   , 0   , 0   }
\definecolor{skyblue}       {cmyk}{0.62, 0   , 0.12, 0   }
\definecolor{turquoise}     {cmyk}{0.85, 0   , 0.20, 0   }
\definecolor{tealblue}      {cmyk}{0.86, 0   , 0.34, 0.02}
\definecolor{aquamarine}    {cmyk}{0.82, 0   , 0.30, 0   }
\definecolor{bluegreen}     {cmyk}{0.85, 0   , 0.33, 0   }
\definecolor{emerald}       {cmyk}{1   , 0   , 0.50, 0   }
\definecolor{junglegreen}   {cmyk}{0.99, 0   , 0.52, 0   }
\definecolor{seagreen}      {cmyk}{0.69, 0   , 0.50, 0   }
\definecolor{green}         {cmyk}{1   , 0   , 1   , 0   }
\definecolor{forestgreen}   {cmyk}{0.91, 0   , 0.88, 0.12}
\definecolor{pinegreen}     {cmyk}{0.92, 0   , 0.59, 0.25}
\definecolor{limegreen}     {cmyk}{0.50, 0   , 1   , 0   }
\definecolor{yellowgreen}   {cmyk}{0.44, 0   , 0.74, 0   }
\definecolor{springgreen}   {cmyk}{0.26, 0   , 0.76, 0   }
\definecolor{olivegreen}    {cmyk}{0.64, 0   , 0.95, 0.40}
\definecolor{rawsienna}     {cmyk}{0   , 0.72, 1   , 0.45}
\definecolor{sepia}         {cmyk}{0   , 0.83, 1   , 0.70}
\definecolor{brown}         {cmyk}{0   , 0.81, 1   , 0.60}
\definecolor{tan}           {cmyk}{0.14, 0.42, 0.56, 0   }
\definecolor{gray}          {cmyk}{0   , 0   , 0   , 0.50}
\definecolor{black}         {cmyk}{0   , 0   , 0   , 1   }
\definecolor{white}         {cmyk}{0   , 0   , 0   , 0   } 

% *************** Koniec definicji stylu dokumentu ***************


% *************** Pakiety dodatkowe ***************
\usepackage{float}
\usepackage{listings}
\usepackage{xcolor}
\usepackage[hidelinks]{hyperref}

\lstset{
	basicstyle=\ttfamily\footnotesize,
	frame=single,
	breaklines=true,
	showstringspaces=false,
	backgroundcolor=\color{black!5},
	literate=
		{ą}{{\k{a}}}1 {ć}{{\'c}}1 {ę}{{\k{e}}}1 {ł}{{\l{}}}1 {ń}{{\'n}}1
		{ó}{{\'o}}1 {ś}{{\'s}}1 {ź}{{\'z}}1 {ż}{{\.z}}1
		{Ą}{{\k{A}}}1 {Ć}{{\'C}}1 {Ę}{{\k{E}}}1 {Ł}{{\L{}}}1 {Ń}{{\'N}}1
		{Ó}{{\'O}}1 {Ś}{{\'S}}1 {Ź}{{\'Z}}1 {Ż}{{\.Z}}1
}

% Zrzut ekranu: \zrzut{plik}{podpis}
\newcommand{\zrzut}[2]{%
	\begin{figure}[H]
		\centering
		\caption{\footnotesize #2}
		\label{fig:#1}
		\includegraphics[width=15.5cm]{figures/#1.png}
	\end{figure}%
}

% *************** Dane do strony tytułowej ***************
\newcommand{\wydzial}{KOLEGIUM INFORMATYKI STOSOWANEJ}
\newcommand{\kierunek}{Kierunek: INFORMATYKA}
\newcommand{\specjalnosc}{}
\newcommand{\autor}{Adam Techman}
\newcommand{\album}{Nr albumu studenta 73285}
\newcommand{\nauczyciel}{mgr inż.~Adrian Hałys}
\newcommand{\tematprojektu}{Rejestr urządzeń sieciowych}
\newcommand{\miasto}{Rzeszów}
\newcommand{\przedmiot}{Programowanie obiektowe}
\newcommand{\rok}{2026}
\newcommand{\repozytorium}{https://github.com/BobnikTheIII/rejestr-urzadzen-sieciowych}

\begin{document}

% *************** Strona tytułowa i spis treści ***************
% *************** Strony tytułowe ***************

% ************************************************************
% W tym miejscu znajduje sie definicja wyglądu pierwszych stron:
% strony tytułowej, strony z oświadczeniem o treści pracy
% i strony ze spisem treści
% ************************************************************
% *************** Strona tytułowa ***************
%umieszczenie logo i nazwy uczelni
\noindent
\parbox{65mm}{\includegraphics[width=13.0cm, height=3.0cm]{logoWSIiZ}}

\vspace{10mm}
\begin{center}
{\Large{}\textbf{\wydzial}}
\end{center}
\vspace{10mm}
\noindent
\hspace{30mm}{\Large{}\textbf{\kierunek}}\\

\noindent
\hspace{30mm}{\Large{}\textbf{\specjalnosc}}
\vspace{30mm}
\begin{center}
	{\large{}\autor}\\
	{\large{}\album}\\
	\vspace{15pt}
	{\huge{}\textbf{Projekt z przedmiotu ,,\przedmiot''}}\\
	\vspace{20pt}
	{\normalsize{}Prowadzący: \nauczyciel}\\
	\vspace{100pt}
	{\LARGE{}\textbf{\tematprojektu}}\\
	\vspace{190pt}
	{\large{}\textbf{\miasto {} \rok}}
\end{center}

% pusta zawartość stopki - brak numeru strony
\thispagestyle{empty}

% *************** Strona z oświadczeniem o treści pracy ***************
\newpage
\text{}

\thispagestyle{empty}
\newpage


% *************** Spis treści ***************
\tableofcontents
% pusta zawartość stopki - brak numeru strony
\thispagestyle{empty}
\newpage

% *************** Koniec pliku front.tex ***************


% *************** Część główna ***************
\chapter*{Wstęp}

Niniejszy dokument opisuje program \textit{Rejestr urządzeń sieciowych}, przygotowany jako projekt z~przedmiotu ,,Programowanie obiektowe''. Jest to konsolowa aplikacja napisana w~języku C\#, realizująca operacje CRUD (Create, Read, Update, Delete) na ewidencji sprzętu sieciowego. Dane są zapisywane w~pliku na dysku, dlatego zachowują się między uruchomieniami programu.

Rozdział \ref{ch:opis} zawiera ogólny opis programu. Rozdział \ref{ch:narzedzia} wymienia użyte narzędzia i~podaje odnośnik do repozytorium z~kodem źródłowym. Rozdział \ref{ch:instrukcja} jest instrukcją obsługi: opisuje szczegółowo każdą funkcję programu i~potwierdza jej działanie zrzutami ekranu. Rozdział \ref{ch:budowa} przedstawia obiektową budowę programu.

\addcontentsline{toc}{chapter}{Wstęp}
\newpage
% ********** Rozdział 1 **********
\chapter{Ogólny opis programu}\label{ch:opis}

\section{Przeznaczenie programu}

Program służy do prowadzenia ewidencji urządzeń sieciowych, takich jak routery, przełączniki, punkty dostępowe i~serwery. Dla każdego urządzenia przechowuje jego nazwę, adres IP, lokalizację i~status, a~także parametry właściwe dla danego typu sprzętu. Może z~niego korzystać na przykład administrator małej sieci, który chce w~jednym miejscu sprawdzić, jaki sprzęt posiada, gdzie się on znajduje i~jakie adresy są już zajęte.

\section{Cel powstania}

Program powstał jako projekt zaliczeniowy z~przedmiotu ,,Programowanie obiektowe''. Jego celem jest praktyczne zastosowanie dziedziczenia, polimorfizmu, hermetyzacji i~abstrakcji w~aplikacji wykonującej trwałe operacje CRUD. Tematyka sieci komputerowych dobrze się do tego nadaje: różne typy urządzeń mają część cech wspólnych i~część odmiennych, co naturalnie prowadzi do hierarchii klas.

\section{Funkcje programu}

Program udostępnia następujące funkcje, wybierane z~menu głównego:

\begin{enumerate}
	\item \textbf{Dodaj urządzenie} (Create) -- wprowadzenie nowego urządzenia wybranego typu.
	\item \textbf{Lista urządzeń} (Read) -- wyświetlenie całego rejestru w~formie tabeli.
	\item \textbf{Szczegóły urządzenia} (Read) -- wyświetlenie wszystkich pól wskazanego urządzenia.
	\item \textbf{Wyszukaj urządzenia} (Read) -- wyszukiwanie tekstu we wszystkich polach urządzeń.
	\item \textbf{Edytuj urządzenie} (Update) -- zmiana wybranych pól istniejącego urządzenia.
	\item \textbf{Usuń urządzenie} (Delete) -- usunięcie urządzenia po potwierdzeniu.
	\item \textbf{Podsumowanie rejestru} -- liczba urządzeń ogółem oraz według typu i~statusu.
\end{enumerate}

Każda zmiana jest zapisywana na dysku natychmiast po jej zatwierdzeniu, bez potrzeby osobnego polecenia zapisu.

\section{Przechowywane dane}

Tabela \ref{tab:pola} przedstawia pola urządzeń. Pola wspólne występują w~każdym typie urządzenia, a~pozostałe zależą od wybranego typu.

\begin{table}[H]
	\begin{center}
		\caption{\footnotesize Pola urządzeń i~reguły ich poprawności}
		\label{tab:pola}
		\begin{small}
			\begin{tabular}{|l|l|l|} \hline
				\textbf{Typ} & \textbf{Pole} & \textbf{Dopuszczalne wartości} \\ \hline
				Wszystkie & Nazwa & 1--40 znaków, unikalna w~rejestrze \\ \cline{2-3}
				& Adres IP & poprawny adres IPv4 lub IPv6, unikalny w~rejestrze \\ \cline{2-3}
				& Lokalizacja & 1--40 znaków \\ \cline{2-3}
				& Status & Aktywne, W~magazynie, W~serwisie, Wycofane \\ \cline{2-3}
				& Uwagi & opcjonalne, do 200 znaków \\ \hline
				Router & Porty WAN & liczba 1--64 \\ \cline{2-3}
				& Obsługa VPN & tak lub nie \\ \hline
				Przełącznik & Liczba portów & liczba 1--256 \\ \cline{2-3}
				& Zarządzalny & tak lub nie \\ \hline
				Punkt dostępowy & SSID & 1--32 znaki \\ \cline{2-3}
				& Pasmo & 2,4 GHz; 5 GHz; 6 GHz; 2,4 + 5 GHz; 2,4 + 5 + 6 GHz \\ \hline
				Serwer & Rdzenie CPU & liczba 1--1024 \\ \cline{2-3}
				& RAM [GB] & liczba 1--65536 \\ \cline{2-3}
				& System operacyjny & 1--40 znaków \\ \hline
			\end{tabular}
		\end{small}
	\end{center}
\end{table}

% ********** Koniec rozdziału **********

\newpage
% ********** Rozdział 2 **********
\chapter{Użyte narzędzia i~kod źródłowy}\label{ch:narzedzia}

\section{Wykaz narzędzi}

Program korzysta wyłącznie z~platformy .NET i~jej biblioteki standardowej. Nie wymaga żadnych dodatkowych pakietów NuGet ani serwera bazy danych. Tabela \ref{tab:narzedzia} zawiera pełny wykaz narzędzi.

\begin{table}[H]
	\begin{center}
		\caption{\footnotesize Wykaz użytych narzędzi}
		\label{tab:narzedzia}
		\begin{small}
			\begin{tabular}{|l|p{10.5cm}|} \hline
				\textbf{Narzędzie} & \textbf{Zastosowanie} \\ \hline
				.NET 9 SDK, język C\# & Kompilacja i~uruchamianie programu. \\ \hline
				System.Text.Json & Zapis i~odczyt danych w~formacie JSON; biblioteka wchodzi w~skład platformy .NET. \\ \hline
				Plik JSON & Trwałe przechowywanie danych; baza danych nie jest używana. \\ \hline
				Git, GitHub & Kontrola wersji i~udostępnienie kodu źródłowego. \\ \hline
				LaTeX (pdfTeX) & Przygotowanie tej dokumentacji na podstawie szablonu uczelnianego. \\ \hline
				Python (Pillow, numpy), Google Chrome & Automatyczne wykonanie zrzutów ekranu do dokumentacji (skrypt \texttt{dokumentacja/narzedzia/zrzuty.py}, opis w~pliku \texttt{ZRZUTY.md}). Nie są potrzebne do działania programu. \\ \hline
			\end{tabular}
		\end{small}
	\end{center}
\end{table}

\section{Repozytorium z~kodem źródłowym}

Kod źródłowy projektu wraz ze źródłami tej dokumentacji znajduje się w~repozytorium:

\begin{center}
	\url{\repozytorium}
\end{center}

\section{Uruchomienie programu}

Do uruchomienia programu potrzebny jest zainstalowany pakiet .NET 9 SDK. Po pobraniu repozytorium program uruchamia się z~jego głównego katalogu poleceniami:

\begin{lstlisting}
git clone https://github.com/BobnikTheIII/rejestr-urzadzen-sieciowych.git
cd rejestr-urzadzen-sieciowych
dotnet run --project src/RejestrUrzadzen
\end{lstlisting}

Program można też otworzyć w~środowisku Visual Studio za pomocą pliku \texttt{RejestrUrzadzen.sln}.

\subsection{Położenie pliku danych}

Domyślnie dane są zapisywane w~pliku \texttt{urzadzenia.json} w~katalogu danych aplikacji bieżącego użytkownika:

\begin{itemize}
	\item Windows: \texttt{\%LOCALAPPDATA\%\textbackslash RejestrUrzadzen\textbackslash urzadzenia.json},
	\item Linux: \texttt{\textasciitilde/.local/share/RejestrUrzadzen/urzadzenia.json}.
\end{itemize}

Inny plik można wskazać jako argument programu. Jeśli plik nie istnieje, zostanie utworzony przy pierwszym zapisie:

\begin{lstlisting}
dotnet run --project src/RejestrUrzadzen -- moj_rejestr.json
\end{lstlisting}

Ścieżka używanego pliku jest zawsze widoczna w~menu głównym.

% ********** Koniec rozdziału **********

\newpage
% ********** Rozdział 3 **********
\chapter{Instrukcja obsługi}\label{ch:instrukcja}

Zrzuty ekranu w~tym rozdziale wykonano automatycznie podczas pracy z~programem: skrypt \texttt{zrzuty.py} uruchomił skompilowany program w~terminalu o~szerokości 100 znaków, wpisał podane dane i~zapisał zawartość ekranu jako obraz w~wyglądzie emulatora terminala Konsole (schemat kolorów Breeze, czcionka Noto Sans Mono) z~systemu Nobara Linux. Rejestr był na początku pusty, dlatego kolejne zrzuty pokazują skutki wcześniejszych operacji. Program wywoływano ze ścieżką względną \texttt{urzadzenia.json}, stąd taka nazwa pliku w~menu.

\section{Menu główne}

Po uruchomieniu program wyświetla menu główne (rysunek \ref{fig:menu}). Nad listą opcji widoczna jest ścieżka używanego pliku danych. Aby wybrać funkcję, należy wpisać jej numer i~nacisnąć Enter.

\zrzut{menu}{Menu główne programu}

Po wykonaniu funkcji program czeka na naciśnięcie klawisza Enter, a~następnie wraca do menu. Wpisanie numeru spoza listy powoduje wyświetlenie komunikatu o~błędzie (rysunek \ref{fig:menu_blad}).

\zrzut{menu_blad}{Reakcja na wybór nieistniejącej opcji}

\section{Dodawanie urządzenia}

Funkcja \textbf{1. Dodaj urządzenie} tworzy nowy wpis w~rejestrze. Należy kolejno:

\begin{enumerate}
	\item wybrać typ urządzenia, podając jego numer,
	\item podać wartości kolejnych pól; w~nawiasach kwadratowych program podpowiada dopuszczalne wartości,
	\item przy polu \textit{Status} oraz \textit{Pasmo} podać numer z~listy lub pełną nazwę wartości,
	\item przy polach typu tak/nie wpisać \texttt{t} lub \texttt{n},
	\item przy opcjonalnym polu \textit{Uwagi} wpisać tekst lub nacisnąć Enter, aby je pominąć.
\end{enumerate}

Po podaniu ostatniego pola urządzenie zostaje zapisane w~pliku, a~program pokazuje nadany mu identyfikator. Identyfikatory nadawane są automatycznie i~nie są używane ponownie, nawet po usunięciu urządzenia. Zestaw pól zależy od typu urządzenia, co pokazują rysunki od \ref{fig:dodaj_router} do \ref{fig:dodaj_serwer}.

\zrzut{dodaj_router}{Dodawanie routera}
\zrzut{dodaj_przelacznik}{Dodawanie przełącznika}
\zrzut{dodaj_ap}{Dodawanie punktu dostępowego}
\zrzut{dodaj_serwer}{Dodawanie serwera}

\subsection{Niepoprawne dane i~anulowanie}

Każda wartość jest sprawdzana od razu po wpisaniu. Gdy jest niepoprawna, program wyświetla przyczynę na czerwono i~ponawia pytanie o~to samo pole, nie tracąc wcześniej podanych danych. Program odrzuca między innymi puste pole wymagane, nazwę lub adres IP zajęty przez inne urządzenie, niepełny adres IP oraz liczbę spoza zakresu. W~każdym momencie można wpisać znak \texttt{!}, aby anulować operację. Urządzenie nie zostanie wtedy dodane. Rysunek \ref{fig:dodaj_anuluj} przedstawia oba te zachowania.

\zrzut{dodaj_anuluj}{Odrzucanie niepoprawnych wartości i~anulowanie dodawania}

\section{Lista urządzeń}

Funkcja \textbf{2. Lista urządzeń} wyświetla wszystkie urządzenia w~kolejności identyfikatorów, wraz z~ich liczbą (rysunek \ref{fig:lista}). Wartości zbyt długie dla szerokości okna konsoli są skracane i~kończą się znakiem wielokropka; pełną wartość pokazuje funkcja szczegółów. Gdy rejestr jest pusty, program informuje o~tym zamiast wyświetlać tabelę (rysunek \ref{fig:lista_pusta}).

\zrzut{lista}{Lista urządzeń}
\zrzut{lista_pusta}{Lista przy pustym rejestrze}

\section{Szczegóły urządzenia}

Funkcja \textbf{3. Szczegóły urządzenia} pokazuje wszystkie pola jednego urządzenia, w~tym pola właściwe dla jego typu. Program wyświetla listę urządzeń i~prosi o~podanie identyfikatora. Jeśli podany tekst nie jest liczbą lub urządzenie o~takim identyfikatorze nie istnieje, program zgłasza błąd i~pyta ponownie. Znak \texttt{!} pozwala wrócić do menu. Niewypełnione pole jest oznaczone znakiem \texttt{-} (rysunek \ref{fig:szczegoly}).

\zrzut{szczegoly}{Szczegóły urządzenia, poprzedzone dwiema niepoprawnymi próbami podania identyfikatora}

\section{Wyszukiwanie urządzeń}

Funkcja \textbf{4. Wyszukaj urządzenia} wyświetla urządzenia, których typ lub dowolne pole zawiera podany tekst. Wielkość liter nie ma znaczenia. Przeszukiwane są także pola właściwe dla danego typu, na przykład SSID lub system operacyjny. Rysunek \ref{fig:szukaj_lokalizacja} przedstawia wyszukiwanie po lokalizacji, a~rysunek \ref{fig:szukaj_pole_typu} -- po systemie operacyjnym serwera. Gdy nic nie pasuje, program wyświetla stosowny komunikat (rysunek \ref{fig:szukaj_brak}).

\zrzut{szukaj_lokalizacja}{Wyszukiwanie po lokalizacji}
\zrzut{szukaj_pole_typu}{Wyszukiwanie po polu właściwym dla typu urządzenia}
\zrzut{szukaj_brak}{Wyszukiwanie bez wyników}

\section{Edycja urządzenia}

Funkcja \textbf{5. Edytuj urządzenie} pozwala zmienić pola istniejącego urządzenia. Po wskazaniu identyfikatora program pyta kolejno o~każde pole, pokazując jego obecną wartość. Dla każdego pola można:

\begin{itemize}
	\item wpisać nową wartość -- obowiązują te same reguły co przy dodawaniu,
	\item nacisnąć Enter, aby zachować obecną wartość,
	\item wpisać \texttt{-} przy polu \textit{Uwagi}, aby je wyczyścić,
	\item wpisać \texttt{!}, aby anulować edycję całego urządzenia.
\end{itemize}

Zmiany są zapisywane dopiero po przejściu przez wszystkie pola. Rysunek \ref{fig:edycja} przedstawia zmianę lokalizacji, statusu, SSID oraz uwag urządzenia nr 3, a~rysunek \ref{fig:edycja_wynik} -- jego stan po zapisaniu.

\zrzut{edycja}{Edycja urządzenia}
\zrzut{edycja_wynik}{Szczegóły urządzenia po edycji}

Typu urządzenia ani jego identyfikatora nie można zmienić. Nazwa i~adres IP nadal muszą być unikalne, przy czym urządzenie może zachować własne wartości. Rysunek \ref{fig:edycja_konflikt} przedstawia próbę użycia nazwy innego urządzenia, a~następnie anulowanie edycji.

\zrzut{edycja_konflikt}{Odrzucenie zajętej nazwy i~anulowanie edycji}

\section{Usuwanie urządzenia}

Funkcja \textbf{6. Usuń urządzenie} trwale usuwa urządzenie z~rejestru. Po wskazaniu identyfikatora program pokazuje dane urządzenia i~prosi o~potwierdzenie. Odpowiedź \texttt{t} usuwa je (rysunek \ref{fig:usun}), a~odpowiedź \texttt{n} pozostawia rejestr bez zmian (rysunek \ref{fig:usun_rezygnacja}).

\zrzut{usun_rezygnacja}{Rezygnacja z~usunięcia urządzenia}
\zrzut{usun}{Usunięcie urządzenia po potwierdzeniu}

\section{Podsumowanie rejestru}

Funkcja \textbf{7. Podsumowanie rejestru} podaje łączną liczbę urządzeń oraz ich liczbę według typu i~statusu (rysunek \ref{fig:podsumowanie}). Rejestr zawiera tu cztery urządzenia, ponieważ jedno z~pięciu dodanych zostało wcześniej usunięte.

\zrzut{podsumowanie}{Podsumowanie rejestru}

\section{Zakończenie pracy}

Opcja \textbf{0. Wyjście} kończy działanie programu (rysunek \ref{fig:wyjscie}). Dane są już wtedy zapisane, ponieważ każda zmiana trafia do pliku od razu.

\zrzut{wyjscie}{Zakończenie pracy programu}

\section{Trwałość danych}

Rysunek \ref{fig:trwalosc} przedstawia listę urządzeń zaraz po ponownym uruchomieniu programu. Widoczne są dane po wcześniejszej edycji i~usunięciu, co potwierdza, że zmiany przetrwały zakończenie procesu.

\zrzut{trwalosc}{Lista urządzeń po ponownym uruchomieniu programu}

Rysunek \ref{fig:trwalosc_id} pokazuje stan po dodaniu kolejnego urządzenia i~następnym uruchomieniu. Nowe urządzenie otrzymało identyfikator 6, a~nie 5, ponieważ identyfikator usuniętego urządzenia nie jest używany ponownie.

\zrzut{trwalosc_id}{Identyfikator usuniętego urządzenia nie wraca do użycia}

\subsection{Format pliku danych}

Dane są przechowywane w~czytelnym pliku JSON. Pole \texttt{type} określa typ urządzenia, a~\texttt{nextId} -- identyfikator kolejnego urządzenia. Listing \ref{lst:json} przedstawia początek pliku z~opisanej sesji.

\noindent
\begin{minipage}{\linewidth}
\begin{lstlisting}[caption={Początek pliku danych},label={lst:json}]
{
  "nextId": 6,
  "devices": [
    {
      "type": "router",
      "wanPorts": 2,
      "supportsVpn": true,
      "id": 1,
      "name": "rtr-brzeg-01",
      "ipAddress": "192.168.0.1",
      "location": "Serwerownia A",
      "status": "Active",
      "notes": "Główny router brzegowy"
    },
\end{lstlisting}
\end{minipage}

\subsection{Uszkodzony plik danych}

Jeśli pliku nie da się poprawnie odczytać, program wyświetla komunikat i~kończy działanie, nie zmieniając pliku (rysunek \ref{fig:plik_uszkodzony}). Dotyczy to między innymi niepoprawnej składni JSON, braku wymaganego pola, wartości niezgodnej z~regułami z~tabeli \ref{tab:pola} oraz powtórzonych identyfikatorów. Plik można wtedy naprawić ręcznie lub wskazać inny jako argument programu. Jeśli natomiast nie uda się zapisać zmiany, na przykład z~powodu braku uprawnień, program wyświetla błąd i~pracuje dalej, a~dane w~pliku pozostają w~poprzednim stanie.

\zrzut{plik_uszkodzony}{Komunikat o~uszkodzonym pliku danych}

% ********** Koniec rozdziału **********

\newpage
% ********** Rozdział 4 **********
\chapter{Budowa programu}\label{ch:budowa}

\section{Podział na warstwy}

Cały kod znajduje się w~klasach. Metoda \texttt{Main} jedynie tworzy obiekty programu i~uruchamia główną pętlę menu. Program składa się z~czterech przestrzeni nazw, przedstawionych w~tabeli \ref{tab:warstwy}.

\begin{table}[H]
	\begin{center}
		\caption{\footnotesize Przestrzenie nazw programu}
		\label{tab:warstwy}
		\begin{small}
			\begin{tabular}{|l|p{11.5cm}|} \hline
				\textbf{Przestrzeń nazw} & \textbf{Odpowiedzialność} \\ \hline
				\texttt{Domain} & Model urządzeń i~reguły poprawności pojedynczych pól. \\ \hline
				\texttt{Persistence} & Trwały zapis i~odczyt danych w~pliku JSON. \\ \hline
				\texttt{Services} & Operacje CRUD, reguła unikalności nazw i~adresów IP, podsumowanie. \\ \hline
				\texttt{Ui} & Komunikacja z~użytkownikiem: menu, formularz, tabela i~polecenia menu. \\ \hline
			\end{tabular}
		\end{small}
	\end{center}
\end{table}

\section{Zastosowanie zasad programowania obiektowego}

\begin{description}
	\item[Dziedziczenie.] Abstrakcyjna klasa \texttt{Device} zawiera cechy wspólne, a~klasy \texttt{Router}, \texttt{NetworkSwitch}, \texttt{AccessPoint} i~\texttt{Server} dodają pola właściwe dla swojego typu. Podobnie po klasie \texttt{MenuCommand} dziedziczy każde polecenie menu, a~po \texttt{DeviceField} -- każdy rodzaj pola.
	\item[Polimorfizm.] Każdy typ urządzenia nadpisuje metodę \texttt{GetSpecificFields}, zwracającą jego pola jako obiekty \texttt{DeviceField}. Formularz, widok szczegółów i~wyszukiwarka korzystają wyłącznie z~tej abstrakcji, dzięki czemu nie zawierają instrukcji warunkowych zależnych od typu urządzenia. Menu główne wywołuje natomiast metodę \texttt{Execute} na liście obiektów \texttt{MenuCommand}.
	\item[Hermetyzacja.] Właściwości modelu sprawdzają wartości w~swoich metodach dostępowych, więc obiekt nie może przyjąć wartości niezgodnej z~regułami, także przy odczycie z~pliku. Identyfikator nadaje wyłącznie warstwa zapisu danych. Repozytorium zwraca kopie obiektów, dlatego zmiany trafiają do rejestru dopiero po jawnym zapisie -- na tym opiera się anulowanie edycji.
	\item[Abstrakcja.] Interfejs \texttt{IDeviceRepository} oddziela logikę programu od sposobu przechowywania danych. Zastąpienie pliku JSON bazą danych wymagałoby tylko nowej klasy implementującej ten interfejs oraz zmiany jednego wiersza w~klasie \texttt{Program}.
\end{description}

\section{Bezpieczeństwo zapisu}

Klasa \texttt{JsonDeviceRepository} zapisuje dane najpierw do pliku tymczasowego, a~następnie zastępuje nim właściwy plik. Przerwanie programu podczas zapisu nie uszkadza więc istniejących danych. Stan w~pamięci zmienia się dopiero po udanym zapisie, dlatego program nie pokazuje zmian, które nie trafiły na dysk.

% ********** Koniec rozdziału **********

\end{document}
% *************** Koniec pliku dokumentacja.tex ***************
